% CV – Rodrigo Blanco Betes
\documentclass[11pt, a4paper]{article}

% ── Paquetes ────────────────────────────────────────────────────────────────
\usepackage[T1]{fontenc}
\usepackage[utf8]{inputenc}
\usepackage[spanish]{babel}
\usepackage[top=2cm, bottom=2cm, left=2cm, right=2cm]{geometry}
\usepackage{lmodern}
\usepackage{microtype}
\usepackage{parskip}
\usepackage{titlesec}
\usepackage{enumitem}
\usepackage{tabularx}
\usepackage{array}
\usepackage{graphicx}
\usepackage{wrapfig}
\usepackage{xcolor}
\usepackage{hyperref}
\usepackage{fontawesome5}

% ── Colores ─────────────────────────────────────────────────────────────────
\definecolor{accentcolor}{RGB}{30, 80, 140}

% ── Hyperref ────────────────────────────────────────────────────────────────
\hypersetup{
  colorlinks=true,
  urlcolor=accentcolor,
  linkcolor=accentcolor,
  pdfauthor={Rodrigo Blanco Betes},
  pdftitle={Curriculum Vitae – Rodrigo Blanco Betes}
}

% ── Secciones ───────────────────────────────────────────────────────────────
\titleformat{\section}
  {\large\bfseries\color{accentcolor}}
  {}{0em}{}
  [\color{accentcolor}\titlerule]
\titlespacing{\section}{0pt}{10pt}{6pt}

% ── Comandos auxiliares ──────────────────────────────────────────────────────
% Entrada de experiencia: {Título}{Fecha}{Institución}{Lugar}
\newcommand{\cventry}[4]{%
  \noindent
  \begin{tabularx}{\linewidth}{@{}X r@{}}
    \textbf{#1} & \textbf{#2}\\
    \textit{#3} & \textit{#4}
  \end{tabularx}%
}

\pagestyle{empty}

% ════════════════════════════════════════════════════════════════════════════
\begin{document}

% ── Cabecera ─────────────────────────────────────────────────────────────────
\begin{wrapfigure}{r}{2.8cm}
  \vspace{-10pt}
%  \includegraphics[width=2.6cm]{photo.jpeg}
\end{wrapfigure}

{\LARGE\bfseries Rodrigo Blanco Betes}

\medskip
\faPhone\ +34 640 344 746 \\
\faEnvelope\ \href{mailto:rodrigoblancobetes@gmail.com}{rodrigoblancobetes@gmail.com} \\
\faMapMarker*\ Bilbao, España

\medskip
Matemático, actualmente doctorando en lógica matemática en la UPV/EHU,
con experiencia en modelización estadística, análisis de datos y desarrollo de
soluciones analíticas en entornos profesionales. Mi formación combina una base
teórica sólida en áreas como análisis y fundamentos matemáticos con experiencia
aplicada en modelización estadística y análisis de datos. Interesado en la
enseñanza de matemáticas con enfoque práctico orientado a la resolución de problemas.

% ── Educación ────────────────────────────────────────────────────────────────
\section{Educación}

\cventry{Doctorado en Lógica y Fundamentos de las Matemáticas}%
        {Oct.\ 2023 -- actualidad}%
        {Universidad del País Vasco, Instituto ILCLI (UPV/EHU)}%
        {San Sebastián, España}
\begin{itemize}[nosep, leftmargin=1.2em, topsep=4pt]
  \item Directores de tesis: María de Ponte Azcárate y Kepa Korta.
  \item Título provisional: \textit{On the undefiniteness of the Continuum Hypothesis
        from a pragmatic linguistic perspective}.
  \item Línea: lógica matemática y fundamentos de las matemáticas.
  \item Beca de investigación predoctoral, Gobierno Vasco (2024--2025).
\end{itemize}

\medskip
\cventry{Grado en Matemáticas}%
        {Sep.\ 2017 -- Jul.\ 2021}%
        {Universidad del País Vasco (UPV/EHU)}%
        {Leioa, España}
\begin{itemize}[nosep, leftmargin=1.2em, topsep=4pt]
  \item Mención en Matemáticas Puras.
  \item Nota media del expediente: 9{,}33/10.
  \item Premio extraordinario y mejor expediente de la promoción.
  \item Programa Erasmus en Lyon, Francia: ENS-Lyon y UCB-Lyon.
\end{itemize}

\medskip
\cventry{Máster en Filosofía: Rama de Lógica y Filosofía de la Ciencia}%
        {Sep.\ 2022 -- Oct.\ 2023}%
        {Universidad Nacional de Educación a Distancia (UNED)}%
        {España}
\begin{itemize}[nosep, leftmargin=1.2em, topsep=4pt]
  \item Nota media del expediente: 9{,}00/10.
\end{itemize}

% ── Docencia y méritos académicos ─────────────────────────────────────────────────────────────────
\section{Docencia y méritos académicos}

\cventry{Docente}%
        {Ene.\ 2026 -- Mar.\ 2026}%
        {Alfa Formación y Consultoría}%
        {Bilbao, España}
\begin{itemize}[nosep, leftmargin=1.2em, topsep=4pt]

  \item Impartición del curso \textit{Especialista en Big Data} (230 horas).
  \item Docencia en las herramientas/lenguajes: Microsoft Power BI, Docker,
        Java, Hadoop, Scala, Apache Spark, Azure.
  \item Aplicación de fundamentos matemáticos (álgebra, estadística, optimización)
        en contextos computacionales como criptografía, programación funcional y modelización.
\end{itemize}

\cventry{Docente}%
        {Sep.\ 2025 -- Oct.\ 2025}%
        {Alfa Formación y Consultoría}%
        {Bilbao, España}
\begin{itemize}[nosep, leftmargin=1.2em, topsep=4pt]
  \item Impartición del curso de formación profesional para el empleo
        \textit{Introducción al Big Data} (12 horas).
  \item Docencia en las herramientas/lenguajes: Microsoft Power BI, Google Locker Studio.
\end{itemize}

\medskip
\cventry{Formación para la gestión y la investigación en la sanidad}%
        {Ene.\ 2024}%
        {Hospital de Santa Marina}%
        {Bilbao, España}
\begin{itemize}[nosep, leftmargin=1.2em, topsep=4pt]
  \item Curso en teoría de juegos y sesgos cognitivos y sus aplicaciones en la sanidad
        (12 horas).
  \item Curso en estadística y análisis de datos para la investigación y gestión
        sanitaria (4 horas).
\end{itemize}

\medskip
\cventry{Seminario (ILCLI)}%
        {Ene.\ 2026}%
        {\textit{An alternative approach to conceptual structuralism}}%
        {Donostia-San Sebastián, España}
\begin{itemize}[nosep, leftmargin=1.2em, topsep=4pt]
  \item Impartición de un seminario de investigación en el marco del ILCLI (Institute for Logic, Cognition, Language and Information).
  \item Desarrollo de una propuesta alternativa al estructuralismo conceptual, abordando sus fundamentos lógicos y filosóficos.
  \item Discusión de implicaciones en filosofía del lenguaje, teoría de conceptos y modelización formal.
  \item Fecha: 30 de enero de 2026.
\end{itemize}

% ── Experiencia profesional ───────────────────────────────────────────────────
\section{Experiencia profesional}

\cventry{Consultor}%
        {Sept.\ 2021 -- Dic.\ 2022}%
        {Management Solutions SL}%
        {Bilbao, España}
\begin{itemize}[nosep, leftmargin=1.2em, topsep=4pt]
  \item Desarrollo y validación de modelos de análisis de datos para riesgo
        de crédito para una entidad bancaria líder en España.
  \item Implementación y análisis de modelos estadísticos para evaluación de riesgo.
  \item Formación interna a nuevos consultores en herramientas de análisis de datos.
  \item Formación adicional en economía y gestión de riesgos en colaboración con
        Universidad Pontificia Comillas -- ICADE.
\end{itemize}

% ── Proyectos de innovación ───────────────────────────────────────────────────
\section{Proyectos de innovación}

\cventry{Co-desarrollador}%
        {Ene.\ 2025 -- Sept.\ 2025}%
        {}{}
\begin{itemize}[nosep, leftmargin=1.2em, topsep=4pt]
  \item Diseño de un sistema para la generación automática de secuencias de turnos
        en entornos sanitarios.
  \item Implementación de criterios de optimización (descansos, carga horaria,
        restricciones laborales).
  \item Desarrollo de prototipo orientado a la mejora de la eficiencia organizativa
        en centros sanitarios.
\end{itemize}

% ── Aptitudes ────────────────────────────────────────────────────────────────
\section{Aptitudes}

\begin{itemize}[nosep, leftmargin=1.2em]
  \item \textbf{Programación:} Python, Scala, Java, SQL.
  \item \textbf{Data \& Big Data:} Apache Spark, Hadoop, Power BI, Azure.
  \item \textbf{Software matemático y estadístico:} MATLAB, Mathematica, R, SPSS.
\end{itemize}

% ── Idiomas ──────────────────────────────────────────────────────────────────
\section{Idiomas}

\begin{tabularx}{\linewidth}{@{}l l l l@{}}
  \textbf{Castellano:} Nativo \\

  \textbf{Inglés:} C1  \\

  \textbf{Euskera:} B1 \\

  \textbf{Alemán:} B2 \\

  \textbf{Francés:} B2 & & &
\end{tabularx}

\end{document}